\chapter{Stability Analysis}

\section{Introduction and Stability Definitions}
Stability is the most critical requirement of any dynamical system. In control engineering, a system is stable if its state remains close to an equilibrium point under small perturbations. Symmetrically, it is unstable if the states diverge under tiny disturbances.

Consider the autonomous LTI continuous-time system:
\begin{equation}
    \dot{x}(t) = A x(t), \quad x(0) = x_0 , \label{eq:stable_ct}
\end{equation}
We define the origin $x=0$ as the equilibrium point, satisfying $A(0) = 0$. If an equilibrium point of interest $x_e$ is non-zero, we can always shift the coordinate system to place the equilibrium at the origin by defining $w = x - x_e$, yielding the homogenous form $\dot{w} = Aw$.

\begin{mydefinition}{Lyapunov Stability}
The equilibrium point $x = 0$ of system \eqref{eq:stable_ct} is stable (in the sense of Lyapunov) if, for every $\epsilon > 0$, there exists a $\delta(\epsilon) > 0$ such that:
\begin{equation}
\| x(t) \| < \epsilon \quad \text{for all } t \ge 0 ,
\end{equation}
whenever the initial state satisfies:
\begin{equation}
\| x_0 \| < \delta .
\end{equation}
\end{mydefinition}
This definition ensures that states starting close to the origin remain bounded within an $\epsilon$-ball for all future time.

\begin{mydefinition}{Asymptotic Stability}
The equilibrium point $x = 0$ is asymptotically stable if:
\begin{enumerate}
    \item It is stable in the sense of Lyapunov.
    \item There exists an $\eta > 0$ such that the state converges to the origin asymptotically:
    \begin{equation}
\lim_{t \to \infty} x(t) = 0 ,
    \end{equation}
    whenever the initial state satisfies $\| x_0 \| < \eta$.
\end{enumerate}
\end{mydefinition}

\begin{mydefinition}{Exponential Stability}
The equilibrium point $x = 0$ is exponentially stable if there exist constants $\alpha > 0$ (decay rate) and $\beta \ge 1$ such that:
\begin{equation}
\| x(t) \| \le \beta \| x_0 \| e^{-\alpha t} \quad \text{for all } t \ge 0 .
\end{equation}
\end{mydefinition}
For LTI systems, asymptotic stability and exponential stability are mathematically equivalent. Symmetrically, a system is stable if and only if all eigenvalues of the dynamic matrix $A$ lie in the open Left-Half of the Complex Plane (LHP), i.e., $\text{Re}(\lambda_i) < 0$ for all $i=1,\dots,n$. Such matrices are called \textbf{Hurwitz} matrices.

\section{Lyapunov Stability Analysis in Continuous-Time}
Aleksandr Lyapunov published his seminal work in 1892, presenting what is now known as Lyapunov's Direct Method. The core idea is to define a generalized scalar "energy" function of the state space and demonstrate that this energy strictly decreases along the state trajectories.

\subsection{Continuous-Time Lyapunov Equation}
For an LTI system $\dot{x} = Ax$, we propose a quadratic Lyapunov function candidate:
\begin{equation}
    V(x) = x^\trp P x , \label{eq:lyap_candidate}
\end{equation}
where $P \in \R^{n \times n}$ is a symmetric positive definite matrix ($P = P^\trp \succ 0$), meaning $V(x) > 0$ for all $x \neq 0$ and $V(0) = 0$.

Differentiating $V(x(t))$ along the system trajectories yields:
\begin{align}
    \dot{V}(x) &= \dot{x}^\trp P x + x^\trp P \dot{x} \nonumber \\
    &= (Ax + Bu)^\trp P x + x^\trp P (Ax + Bu) \nonumber \\
&= x^\trp (A^\trp P + P A) x, \quad \text{for } u(t) = 0 ,
\end{align}
For asymptotic stability, we require the energy to strictly decrease along all trajectories, i.e., $\dot{V}(x) < 0$ for all $x \neq 0$. This holds if and only if the matrix term inside is negative definite:
\begin{equation}
A^\trp P + P A \prec 0 ,
\end{equation}
Defining this negative definite term as $-Q$ for a chosen symmetric positive definite matrix $Q \succ 0$ yields the continuous-time \textbf{Lyapunov Matrix Equation}:
\begin{equation}
    A^\trp P + P A + Q = 0 , \label{eq:lyap_eq}
\end{equation}

\begin{mytheorem}{Continuous-Time Lyapunov Stability}
The matrix $A$ is Hurwitz (stable) if and only if, for any symmetric positive definite matrix $Q \succ 0$, there exists a unique symmetric positive definite matrix $P \succ 0$ satisfying the Lyapunov equation \eqref{eq:lyap_eq}.
\end{mytheorem}

\begin{proof}
We prove both directions of the theorem:
\begin{enumerate}
    \item \textbf{Sufficiency ($\Leftarrow$)}: Assume there exist $P \succ 0, Q \succ 0$ satisfying \eqref{eq:lyap_eq}. Let $\lambda$ be an eigenvalue of $A$ with corresponding eigenvector $v \neq 0$, so $Av = \lambda v$.
    Taking the conjugate transpose yields $v^* A^\trp = \lambda^* v^*$.
    Premultiplying \eqref{eq:lyap_eq} by $v^*$ and postmultiplying by $v$ gives:
    \begin{align}
        v^* (A^\trp P + P A + Q) v &= 0 \nonumber \\
        v^* A^\trp P v + v^* P A v + v^* Q v &= 0 \nonumber \\
        \lambda^* v^* P v + \lambda v^* P v + v^* Q v &= 0 \nonumber \\
2 \text{Re}(\lambda) (v^* P v) &= -v^* Q v ,
    \end{align}
    Since $P \succ 0$ and $Q \succ 0$, we have $v^* P v > 0$ and $v^* Q v > 0$. This implies:
    \begin{equation}
\text{Re}(\lambda) = -\frac{v^* Q v}{2 v^* P v} < 0 ,
    \end{equation}
    Thus, all eigenvalues of $A$ have strictly negative real parts, proving stability.
    \item \textbf{Necessity ($\Rightarrow$)}: Assume $A$ is stable. We propose the analytical solution:
    \begin{equation}
P = \int_0^\infty e^{A^\trp \tau} Q e^{A \tau} d\tau ,
    \end{equation}
    Since $A$ is stable, $e^{A\tau} \to 0$ exponentially as $\tau \to \infty$, guaranteeing that the integral converges. Since $Q \succ 0$, it is straightforward that $P \succ 0$.
    To verify that $P$ satisfies the Lyapunov equation, we substitute it:
    \begin{align}
        A^\trp P + P A &= A^\trp \left( \int_0^\infty e^{A^\trp \tau} Q e^{A \tau} d\tau \right) + \left( \int_0^\infty e^{A^\trp \tau} Q e^{A \tau} d\tau \right) A \nonumber \\
        &= \int_0^\infty \left( A^\trp e^{A^\trp \tau} Q e^{A \tau} + e^{A^\trp \tau} Q e^{A \tau} A \right) d\tau \nonumber \\
        &= \int_0^\infty \frac{d}{d\tau} \left( e^{A^\trp \tau} Q e^{A \tau} \right) d\tau \nonumber \\
&= \left. e^{A^\trp \tau} Q e^{A \tau} \right|_0^\infty = 0 - Q = -Q ,
    \end{align}
    This proves that $P$ is indeed the unique positive definite solution.
\end{enumerate}
\end{proof}

\section{Lyapunov Stability Analysis in Discrete-Time}
Consider the unforced discrete-time LTI system:
\begin{equation}
    x_{k+1} = A x_k, \quad x(0) = x_0 , \label{eq:stable_dt}
\end{equation}
The system is asymptotically stable if and only if all eigenvalues of $A$ lie in the open unit circle of the complex plane, i.e., $|\lambda_i| < 1$ for all $i=1,\dots,n$. Such matrices are called \textbf{Schur} matrices.

\subsection{Discrete-Time Lyapunov (Stein) Equation}
Using the quadratic Lyapunov candidate $V(x_k) = x_k^\trp P x_k$ with $P \succ 0$, we require the energy difference between successive samples to be strictly negative:
\begin{align}
    \Delta V(x_k) &= V(x_{k+1}) - V(x_k) \nonumber \\
    &= x_{k+1}^\trp P x_{k+1} - x_k^\trp P x_k \nonumber \\
    &= (A x_k)^\trp P (A x_k) - x_k^\trp P x_k \nonumber \\
&= x_k^\trp (A^\trp P A - P) x_k < 0 ,
\end{align}
Defining this negative definite term as $-Q$ for a chosen symmetric $Q \succ 0$ yields the \textbf{Stein (Discrete Lyapunov) Equation}:
\begin{equation}
    A^\trp P A - P + Q = 0 , \label{eq:stein_eq}
\end{equation}

\begin{mytheorem}{Discrete-Time Lyapunov Stability}
The matrix $A$ is Schur (stable) if and only if, for any symmetric positive definite matrix $Q \succ 0$, there exists a unique symmetric positive definite matrix $P \succ 0$ satisfying the Stein equation \eqref{eq:stein_eq}.
\end{mytheorem}

\section{Polytopic Uncertainty and Quadratic Stability}
In physical systems, matrices $A$ are often not perfectly known due to modeling errors or operating variations. A powerful model for such uncertainty is the \textbf{Polytopic Uncertainty Model}.

Let the dynamic matrix $A(t)$ belong to a polytope of matrices defined by the convex hull of $M$ vertex matrices $\{A_1, A_2, \dots, A_M\}$:
\begin{equation}
A(t) \in \text{Co}\{A_1, \dots, A_M\} = \left\{ \sum_{i=1}^M \alpha_i(t) A_i \ \middle| \ \sum_{i=1}^M \alpha_i(t) = 1, \ \alpha_i(t) \ge 0 \right\} ,
\end{equation}
Evaluating the stability of the time-varying system $\dot{x}(t) = A(t) x(t)$ for infinitely many possible values of $\alpha_i(t)$ is impossible. However, we can establish stability using a \textbf{Common Quadratic Lyapunov Function}.

\begin{mydefinition}{Quadratic Stability}
The uncertain system $\dot{x}(t) = A(t) x(t)$ is quadratically stable if there exists a single, common symmetric positive definite matrix $P \succ 0$ such that:
\begin{equation}
A(t)^\trp P + P A(t) \prec 0 \quad \text{for all possible } A(t) .
\end{equation}
\end{mydefinition}

\begin{mytheorem}{LMI Vertex Stability Condition}
The uncertain system is quadratically stable if and only if there exists a single symmetric matrix $P$ satisfying the system of \textbf{Linear Matrix Inequalities (LMIs)}:
\begin{align}
    P &\succ 0 \\
A_i^\trp P + P A_i &\prec 0, \quad \forall i = 1, \dots, M .
\end{align}
\end{mytheorem}
\begin{proof}
Because $A(t) = \sum_{i=1}^M \alpha_i(t) A_i$ is a convex combination:
\begin{align}
    A(t)^\trp P + P A(t) &= \left( \sum_{i=1}^M \alpha_i(t) A_i \right)^\trp P + P \left( \sum_{i=1}^M \alpha_i(t) A_i \right) \nonumber \\
&= \sum_{i=1}^M \alpha_i(t) \left( A_i^\trp P + P A_i \right) ,
\end{align}
Since $\alpha_i(t) \ge 0$ and $\sum \alpha_i = 1$, if each vertex matrix satisfies the LMI $A_i^\trp P + P A_i \prec 0$, then their convex sum is guaranteed to be strictly negative definite. This reduces an infinite-dimensional stability problem to checking a finite set of $M$ LMIs at the vertices of the polytope, which can be solved numerically in milliseconds using semi-definite programming (SDP) solvers!
\end{proof}
